% Options for packages loaded elsewhere
\PassOptionsToPackage{unicode}{hyperref}
\PassOptionsToPackage{hyphens}{url}
\documentclass[
]{article}
\usepackage{xcolor}
\usepackage{amsmath,amssymb}
\setcounter{secnumdepth}{-\maxdimen} % remove section numbering
\usepackage{iftex}
\ifPDFTeX
  \usepackage[T1]{fontenc}
  \usepackage[utf8]{inputenc}
  \usepackage{textcomp} % provide euro and other symbols
\else % if luatex or xetex
  \usepackage{unicode-math} % this also loads fontspec
  \defaultfontfeatures{Scale=MatchLowercase}
  \defaultfontfeatures[\rmfamily]{Ligatures=TeX,Scale=1}
\fi
\usepackage{lmodern}
\ifPDFTeX\else
  % xetex/luatex font selection
\fi
% Use upquote if available, for straight quotes in verbatim environments
\IfFileExists{upquote.sty}{\usepackage{upquote}}{}
\IfFileExists{microtype.sty}{% use microtype if available
  \usepackage[]{microtype}
  \UseMicrotypeSet[protrusion]{basicmath} % disable protrusion for tt fonts
}{}
\makeatletter
\@ifundefined{KOMAClassName}{% if non-KOMA class
  \IfFileExists{parskip.sty}{%
    \usepackage{parskip}
  }{% else
    \setlength{\parindent}{0pt}
    \setlength{\parskip}{6pt plus 2pt minus 1pt}}
}{% if KOMA class
  \KOMAoptions{parskip=half}}
\makeatother
\usepackage{color}
\usepackage{fancyvrb}
\newcommand{\VerbBar}{|}
\newcommand{\VERB}{\Verb[commandchars=\\\{\}]}
\DefineVerbatimEnvironment{Highlighting}{Verbatim}{commandchars=\\\{\}}
% Add ',fontsize=\small' for more characters per line
\newenvironment{Shaded}{}{}
\newcommand{\AlertTok}[1]{\textcolor[rgb]{1.00,0.00,0.00}{\textbf{#1}}}
\newcommand{\AnnotationTok}[1]{\textcolor[rgb]{0.38,0.63,0.69}{\textbf{\textit{#1}}}}
\newcommand{\AttributeTok}[1]{\textcolor[rgb]{0.49,0.56,0.16}{#1}}
\newcommand{\BaseNTok}[1]{\textcolor[rgb]{0.25,0.63,0.44}{#1}}
\newcommand{\BuiltInTok}[1]{\textcolor[rgb]{0.00,0.50,0.00}{#1}}
\newcommand{\CharTok}[1]{\textcolor[rgb]{0.25,0.44,0.63}{#1}}
\newcommand{\CommentTok}[1]{\textcolor[rgb]{0.38,0.63,0.69}{\textit{#1}}}
\newcommand{\CommentVarTok}[1]{\textcolor[rgb]{0.38,0.63,0.69}{\textbf{\textit{#1}}}}
\newcommand{\ConstantTok}[1]{\textcolor[rgb]{0.53,0.00,0.00}{#1}}
\newcommand{\ControlFlowTok}[1]{\textcolor[rgb]{0.00,0.44,0.13}{\textbf{#1}}}
\newcommand{\DataTypeTok}[1]{\textcolor[rgb]{0.56,0.13,0.00}{#1}}
\newcommand{\DecValTok}[1]{\textcolor[rgb]{0.25,0.63,0.44}{#1}}
\newcommand{\DocumentationTok}[1]{\textcolor[rgb]{0.73,0.13,0.13}{\textit{#1}}}
\newcommand{\ErrorTok}[1]{\textcolor[rgb]{1.00,0.00,0.00}{\textbf{#1}}}
\newcommand{\ExtensionTok}[1]{#1}
\newcommand{\FloatTok}[1]{\textcolor[rgb]{0.25,0.63,0.44}{#1}}
\newcommand{\FunctionTok}[1]{\textcolor[rgb]{0.02,0.16,0.49}{#1}}
\newcommand{\ImportTok}[1]{\textcolor[rgb]{0.00,0.50,0.00}{\textbf{#1}}}
\newcommand{\InformationTok}[1]{\textcolor[rgb]{0.38,0.63,0.69}{\textbf{\textit{#1}}}}
\newcommand{\KeywordTok}[1]{\textcolor[rgb]{0.00,0.44,0.13}{\textbf{#1}}}
\newcommand{\NormalTok}[1]{#1}
\newcommand{\OperatorTok}[1]{\textcolor[rgb]{0.40,0.40,0.40}{#1}}
\newcommand{\OtherTok}[1]{\textcolor[rgb]{0.00,0.44,0.13}{#1}}
\newcommand{\PreprocessorTok}[1]{\textcolor[rgb]{0.74,0.48,0.00}{#1}}
\newcommand{\RegionMarkerTok}[1]{#1}
\newcommand{\SpecialCharTok}[1]{\textcolor[rgb]{0.25,0.44,0.63}{#1}}
\newcommand{\SpecialStringTok}[1]{\textcolor[rgb]{0.73,0.40,0.53}{#1}}
\newcommand{\StringTok}[1]{\textcolor[rgb]{0.25,0.44,0.63}{#1}}
\newcommand{\VariableTok}[1]{\textcolor[rgb]{0.10,0.09,0.49}{#1}}
\newcommand{\VerbatimStringTok}[1]{\textcolor[rgb]{0.25,0.44,0.63}{#1}}
\newcommand{\WarningTok}[1]{\textcolor[rgb]{0.38,0.63,0.69}{\textbf{\textit{#1}}}}
\usepackage{longtable,booktabs,array}
\usepackage{caption}
\captionsetup[table]{skip=6pt}
\newcounter{none} % for unnumbered tables
\usepackage{calc} % for calculating minipage widths
% Correct order of tables after \paragraph or \subparagraph
\usepackage{etoolbox}
\makeatletter
\patchcmd\longtable{\par}{\if@noskipsec\mbox{}\fi\par}{}{}
\makeatother
% Allow footnotes in longtable head/foot
\IfFileExists{footnotehyper.sty}{\usepackage{footnotehyper}}{\usepackage{footnote}}
\makesavenoteenv{longtable}
\usepackage{graphicx}
\makeatletter
\newsavebox\pandoc@box
\newcommand*\pandocbounded[1]{% scales image to fit in text height/width
  \sbox\pandoc@box{#1}%
  \Gscale@div\@tempa{\textheight}{\dimexpr\ht\pandoc@box+\dp\pandoc@box\relax}%
  \Gscale@div\@tempb{\linewidth}{\wd\pandoc@box}%
  \ifdim\@tempb\p@<\@tempa\p@\let\@tempa\@tempb\fi% select the smaller of both
  \ifdim\@tempa\p@<\p@\scalebox{\@tempa}{\usebox\pandoc@box}%
  \else\usebox{\pandoc@box}%
  \fi%
}
% Set default figure placement to htbp
\def\fps@figure{htbp}
\makeatother
\ifLuaTeX
  \usepackage{luacolor}
  \usepackage[soul]{lua-ul}
\else
  \usepackage{soul}
\fi
% definitions for citeproc citations
\NewDocumentCommand\citeproctext{}{}
\NewDocumentCommand\citeproc{mm}{%
  \begingroup\def\citeproctext{#2}\cite{#1}\endgroup}
\makeatletter
 % allow citations to break across lines
 \let\@cite@ofmt\@firstofone
 % avoid brackets around text for \cite:
 \def\@biblabel#1{}
 \def\@cite#1#2{{#1\if@tempswa , #2\fi}}
\makeatother
\newlength{\cslhangindent}
\setlength{\cslhangindent}{1.5em}
\newlength{\csllabelwidth}
\setlength{\csllabelwidth}{3em}
\newenvironment{CSLReferences}[2] % #1 hanging-indent, #2 entry-spacing
 {\begin{list}{}{%
  \setlength{\itemindent}{0pt}
  \setlength{\leftmargin}{0pt}
  \setlength{\parsep}{0pt}
  % turn on hanging indent if param 1 is 1
  \ifodd #1
   \setlength{\leftmargin}{\cslhangindent}
   \setlength{\itemindent}{-1\cslhangindent}
  \fi
  % set entry spacing
  \setlength{\itemsep}{#2\baselineskip}}}
 {\end{list}}
\usepackage{calc}
\newcommand{\CSLBlock}[1]{\hfill\break\parbox[t]{\linewidth}{\strut\ignorespaces#1\strut}}
\newcommand{\CSLLeftMargin}[1]{\parbox[t]{\csllabelwidth}{\strut#1\strut}}
\newcommand{\CSLRightInline}[1]{\parbox[t]{\linewidth - \csllabelwidth}{\strut#1\strut}}
\newcommand{\CSLIndent}[1]{\hspace{\cslhangindent}#1}
\setlength{\emergencystretch}{3em} % prevent overfull lines
\providecommand{\tightlist}{%
  \setlength{\itemsep}{0pt}\setlength{\parskip}{0pt}}
\usepackage{bookmark}
\IfFileExists{xurl.sty}{\usepackage{xurl}}{} % add URL line breaks if available
\urlstyle{same}
% fallback for those not using the hyperref driver hyperxmp:
\makeatletter
\@ifundefined{xmpquote}{\newcommand{\xmpquote}[1]{#1}}{}
\makeatother
\hypersetup{
  hidelinks,
  pdfcreator={LaTeX via pandoc}}

\author{}
\date{}

\begin{document}

\section{Halmos Test Document}\label{halmos-test-document}

A tour of everything Halmos renders, written to be read as a PDF. Every
section shows the markdown first, in a code block, and then the same
text as Halmos renders it --- so the page itself is the proof. Open this
file in Halmos and press ⌘E to export it.

Swap \texttt{csl:} in the frontmatter to \texttt{ieee},
\texttt{vancouver}, \texttt{chicago-author-date} or \texttt{harvard} to
re-render every citation and the References list in that style.

\textbf{About the bibliography.} \texttt{zotero-export-text.bib} is
auto-synced from Zotero --- do not hand-edit it, and do not add fixture
entries to it. The citations section names specific works, so if the
library changes, those expectations need updating rather than the file.

\subsection{1. Text basics}\label{text-basics}

\subsubsection{Markdown}\label{markdown}

\begin{Shaded}
\begin{Highlighting}[]
\NormalTok{This paragraph has **bold text**, *italic text*, ***bold italic***, \textasciigrave{}inline}
\NormalTok{code\textasciigrave{}, \textasciitilde{}\textasciitilde{}strikethrough\textasciitilde{}\textasciitilde{}, and a }\CommentTok{[}\OtherTok{link to the Wails docs}\CommentTok{](https://v3.wails.io)}\NormalTok{.}

\NormalTok{A second paragraph, to show spacing.}

\AttributeTok{\textgreater{} A blockquote: "The purpose of computing is insight, not numbers."}

\NormalTok{{-}{-}{-}}

\NormalTok{That was a horizontal rule.}
\end{Highlighting}
\end{Shaded}

\subsubsection{Rendered}\label{rendered}

This paragraph has \textbf{bold text}, \emph{italic text},
\textbf{\emph{bold italic}}, \texttt{inline\ code}, \st{strikethrough},
and a \href{https://v3.wails.io}{link to the Wails docs}.

A second paragraph, to show spacing.

\begin{quote}
A blockquote: ``The purpose of computing is insight, not numbers.''
\end{quote}

\begin{center}\rule{0.5\linewidth}{0.5pt}\end{center}

That was a horizontal rule.

\subsection{2. Lists}\label{lists}

\subsubsection{Markdown}\label{markdown-1}

\begin{Shaded}
\begin{Highlighting}[]
\SpecialStringTok{1. }\NormalTok{First ordered item}
\SpecialStringTok{2. }\NormalTok{Second ordered item}
\SpecialStringTok{   1. }\NormalTok{Nested ordered child}
\SpecialStringTok{   2. }\NormalTok{Another nested child}
\SpecialStringTok{3. }\NormalTok{Third ordered item}

\SpecialStringTok{{-} }\NormalTok{Unordered item with a sub{-}list:}
\SpecialStringTok{  {-} }\NormalTok{Nested bullet one}
\SpecialStringTok{  {-} }\NormalTok{Nested bullet two}
\SpecialStringTok{{-} }\NormalTok{Unordered item with }\InformationTok{\textasciigrave{}code\textasciigrave{}}\NormalTok{ and **bold** inside}
\end{Highlighting}
\end{Shaded}

\subsubsection{Rendered}\label{rendered-1}

\begin{enumerate}
\def\labelenumi{\arabic{enumi}.}
\tightlist
\item
  First ordered item
\item
  Second ordered item

  \begin{enumerate}
  \def\labelenumii{\arabic{enumii}.}
  \tightlist
  \item
    Nested ordered child
  \item
    Another nested child
  \end{enumerate}
\item
  Third ordered item
\end{enumerate}

\begin{itemize}
\tightlist
\item
  Unordered item with a sub-list:

  \begin{itemize}
  \tightlist
  \item
    Nested bullet one
  \item
    Nested bullet two
  \end{itemize}
\item
  Unordered item with \texttt{code} and \textbf{bold} inside
\end{itemize}

\subsection{3. Table}\label{table}

The \texttt{:-\/-\/-:} marker centres a column, and maths works inside a
cell.

\subsubsection{Markdown}\label{markdown-2}

\begin{Shaded}
\begin{Highlighting}[]
\PreprocessorTok{|}\NormalTok{ Quantity }\PreprocessorTok{|}\NormalTok{ Symbol }\PreprocessorTok{|}\NormalTok{ Unit }\PreprocessorTok{|}
\PreprocessorTok{|{-}{-}{-}{-}{-}{-}{-}{-}{-}{-}|:{-}{-}{-}{-}{-}{-}:|{-}{-}{-}{-}{-}{-}|}
\PreprocessorTok{|}\NormalTok{ Energy }\PreprocessorTok{|}\NormalTok{ $E$ }\PreprocessorTok{|}\NormalTok{ joule }\PreprocessorTok{|}
\PreprocessorTok{|}\NormalTok{ Entropy }\PreprocessorTok{|}\NormalTok{ $S$ }\PreprocessorTok{|}\NormalTok{ J/K }\PreprocessorTok{|}
\PreprocessorTok{|}\NormalTok{ Density matrix }\PreprocessorTok{|}\NormalTok{ $\textbackslash{}rho\_\textbackslash{}mathcal\{S\}$ }\PreprocessorTok{|}\NormalTok{ — }\PreprocessorTok{|}
\end{Highlighting}
\end{Shaded}

\subsubsection{Rendered}\label{rendered-2}

{\def\LTcaptype{none} % do not increment counter
\begin{longtable}[]{@{}lcl@{}}
\toprule\noalign{}
Quantity & Symbol & Unit \\
\midrule\noalign{}
\endhead
\bottomrule\noalign{}
\endlastfoot
Energy & \(E\) & joule \\
Entropy & \(S\) & J/K \\
Density matrix & \(\rho_\mathcal{S}\) & --- \\
\end{longtable}
}

\subsection{4. Code blocks}\label{code-blocks}

A fence naming a language is syntax-coloured; a fence naming a language
Halmos does not know renders as plain text rather than erroring.

\subsubsection{Markdown}\label{markdown-3}

\begin{Shaded}
\begin{Highlighting}[]
\InformationTok{\textasciigrave{}\textasciigrave{}\textasciigrave{}python}
\KeywordTok{def}\NormalTok{ entropy(p):}
    \CommentTok{"""Shannon entropy."""}
    \ControlFlowTok{return} \OperatorTok{{-}}\BuiltInTok{sum}\NormalTok{(x }\OperatorTok{*}\NormalTok{ math.log2(x) }\ControlFlowTok{for}\NormalTok{ x }\KeywordTok{in}\NormalTok{ p }\ControlFlowTok{if}\NormalTok{ x }\OperatorTok{\textgreater{}} \DecValTok{0}\NormalTok{)}
\InformationTok{\textasciigrave{}\textasciigrave{}\textasciigrave{}}

\InformationTok{\textasciigrave{}\textasciigrave{}\textasciigrave{}notalang}
\InformationTok{this fence names a language nobody has heard of}
\InformationTok{\textasciigrave{}\textasciigrave{}\textasciigrave{}}
\end{Highlighting}
\end{Shaded}

\subsubsection{Rendered}\label{rendered-3}

\begin{Shaded}
\begin{Highlighting}[]
\KeywordTok{def}\NormalTok{ entropy(p):}
    \CommentTok{"""Shannon entropy."""}
    \ControlFlowTok{return} \OperatorTok{{-}}\BuiltInTok{sum}\NormalTok{(x }\OperatorTok{*}\NormalTok{ math.log2(x) }\ControlFlowTok{for}\NormalTok{ x }\KeywordTok{in}\NormalTok{ p }\ControlFlowTok{if}\NormalTok{ x }\OperatorTok{\textgreater{}} \DecValTok{0}\NormalTok{)}
\end{Highlighting}
\end{Shaded}

\begin{Shaded}
\begin{Highlighting}[]
\NormalTok{this fence names a language nobody has heard of}
\end{Highlighting}
\end{Shaded}

\subsection{5. Inline maths}\label{inline-maths}

Single dollars delimit inline maths. Plain dollar amounts are left
alone.

\subsubsection{Markdown}\label{markdown-4}

\begin{Shaded}
\begin{Highlighting}[]
\NormalTok{Euler\textquotesingle{}s identity $e\^{}\{i\textbackslash{}pi\} + 1 = 0$ sits inside a sentence. The mass–energy}
\NormalTok{relation $E = mc\^{}2$, a subscripted state $\textbackslash{}rho\_\textbackslash{}mathcal\{S\}$, a fraction}
\NormalTok{$\textbackslash{}tfrac\{a+b\}\{c\}$, and Greek letters $\textbackslash{}alpha, \textbackslash{}beta, \textbackslash{}gamma$ all sit on the}
\NormalTok{surrounding baseline. Plain dollar amounts like $5 and $10 do not become maths.}
\end{Highlighting}
\end{Shaded}

\subsubsection{Rendered}\label{rendered-4}

Euler's identity \(e^{i\pi} + 1 = 0\) sits inside a sentence. The
mass--energy relation \(E = mc^2\), a subscripted state
\(\rho_\mathcal{S}\), a fraction \(\tfrac{a+b}{c}\), and Greek letters
\(\alpha, \beta, \gamma\) all sit on the surrounding baseline. Plain
dollar amounts like \$5 and \$10 do not become maths.

\subsection{6. Display maths}\label{display-maths}

Double dollars set an equation on its own line, centred.

\subsubsection{Markdown}\label{markdown-5}

\begin{Shaded}
\begin{Highlighting}[]
\AnnotationTok{The Gaussian integral:}

\NormalTok{$$\textbackslash{}int\_\{{-}\textbackslash{}infty\}\^{}\{\textbackslash{}infty\} e\^{}\{{-}x\^{}2\}\textbackslash{},dx = \textbackslash{}sqrt\{\textbackslash{}pi\}$$}

\NormalTok{A matrix and a summation:}

\NormalTok{$$A = \textbackslash{}begin\{pmatrix\} 1 \& 2 }\SpecialCharTok{\textbackslash{}\textbackslash{}}\NormalTok{ 3 \& 4 \textbackslash{}end\{pmatrix\}, \textbackslash{}qquad \textbackslash{}sum\_\{n=1\}\^{}\{\textbackslash{}infty\} \textbackslash{}frac\{1\}\{n\^{}2\} = \textbackslash{}frac\{\textbackslash{}pi\^{}2\}\{6\}$$}

\NormalTok{Aligned equations:}

\NormalTok{$$\textbackslash{}begin\{aligned\} \textbackslash{}nabla \textbackslash{}cdot \textbackslash{}mathbf\{E\} \&= \textbackslash{}frac\{\textbackslash{}rho\}\{\textbackslash{}varepsilon\_0\} }\SpecialCharTok{\textbackslash{}\textbackslash{}}\NormalTok{ \textbackslash{}nabla \textbackslash{}cdot \textbackslash{}mathbf\{B\} \&= 0 \textbackslash{}end\{aligned\}$$}
\end{Highlighting}
\end{Shaded}

\subsubsection{Rendered}\label{rendered-5}

The Gaussian integral:

\[\int_{-\infty}^{\infty} e^{-x^2}\,dx = \sqrt{\pi}\]

A matrix and a summation:

\[A = \begin{pmatrix} 1 & 2 \\ 3 & 4 \end{pmatrix}, \qquad \sum_{n=1}^{\infty} \frac{1}{n^2} = \frac{\pi^2}{6}\]

Aligned equations:

\[\begin{aligned} \nabla \cdot \mathbf{E} &= \frac{\rho}{\varepsilon_0} \\ \nabla \cdot \mathbf{B} &= 0 \end{aligned}\]

\subsection{7. Charts}\label{charts}

A \texttt{vega-lite} fence containing a Vega-Lite spec becomes a live
chart. The data can be inlined, as here, or built with Insert →
Chart\ldots{} from a pasted or imported table.

\subsubsection{Markdown}\label{markdown-6}

\begin{Shaded}
\begin{Highlighting}[]
\InformationTok{\textasciigrave{}\textasciigrave{}\textasciigrave{}vega{-}lite}
\InformationTok{\{}
\InformationTok{  "description": "Simple bar chart",}
\InformationTok{  "data": \{}
\InformationTok{    "values": [}
\InformationTok{      \{"category": "A", "value": 28\}, \{"category": "B", "value": 55\},}
\InformationTok{      \{"category": "C", "value": 43\}, \{"category": "D", "value": 91\},}
\InformationTok{      \{"category": "E", "value": 81\}, \{"category": "F", "value": 53\}}
\InformationTok{    ]}
\InformationTok{  \},}
\InformationTok{  "mark": "bar",}
\InformationTok{  "encoding": \{}
\InformationTok{    "x": \{"field": "category", "type": "nominal", "axis": \{"labelAngle": 0\}\},}
\InformationTok{    "y": \{"field": "value", "type": "quantitative"\}}
\InformationTok{  \}}
\InformationTok{\}}
\InformationTok{\textasciigrave{}\textasciigrave{}\textasciigrave{}}

\InformationTok{\textasciigrave{}\textasciigrave{}\textasciigrave{}vega{-}lite}
\InformationTok{\{}
\InformationTok{  "description": "Line chart with points",}
\InformationTok{  "data": \{}
\InformationTok{    "values": [}
\InformationTok{      \{"x": 0, "y": 1.0\}, \{"x": 1, "y": 1.6\}, \{"x": 2, "y": 2.6\},}
\InformationTok{      \{"x": 3, "y": 4.1\}, \{"x": 4, "y": 6.6\}, \{"x": 5, "y": 10.5\}}
\InformationTok{    ]}
\InformationTok{  \},}
\InformationTok{  "mark": \{"type": "line", "point": true\},}
\InformationTok{  "encoding": \{}
\InformationTok{    "x": \{"field": "x", "type": "quantitative"\},}
\InformationTok{    "y": \{"field": "y", "type": "quantitative"\}}
\InformationTok{  \}}
\InformationTok{\}}
\InformationTok{\textasciigrave{}\textasciigrave{}\textasciigrave{}}
\end{Highlighting}
\end{Shaded}

\subsubsection{Rendered}\label{rendered-6}

\begin{Shaded}
\begin{Highlighting}[]
\NormalTok{\{}
\NormalTok{  "description": "Simple bar chart",}
\NormalTok{  "data": \{}
\NormalTok{    "values": [}
\NormalTok{      \{"category": "A", "value": 28\}, \{"category": "B", "value": 55\},}
\NormalTok{      \{"category": "C", "value": 43\}, \{"category": "D", "value": 91\},}
\NormalTok{      \{"category": "E", "value": 81\}, \{"category": "F", "value": 53\}}
\NormalTok{    ]}
\NormalTok{  \},}
\NormalTok{  "mark": "bar",}
\NormalTok{  "encoding": \{}
\NormalTok{    "x": \{"field": "category", "type": "nominal", "axis": \{"labelAngle": 0\}\},}
\NormalTok{    "y": \{"field": "value", "type": "quantitative"\}}
\NormalTok{  \}}
\NormalTok{\}}
\end{Highlighting}
\end{Shaded}

\begin{Shaded}
\begin{Highlighting}[]
\NormalTok{\{}
\NormalTok{  "description": "Line chart with points",}
\NormalTok{  "data": \{}
\NormalTok{    "values": [}
\NormalTok{      \{"x": 0, "y": 1.0\}, \{"x": 1, "y": 1.6\}, \{"x": 2, "y": 2.6\},}
\NormalTok{      \{"x": 3, "y": 4.1\}, \{"x": 4, "y": 6.6\}, \{"x": 5, "y": 10.5\}}
\NormalTok{    ]}
\NormalTok{  \},}
\NormalTok{  "mark": \{"type": "line", "point": true\},}
\NormalTok{  "encoding": \{}
\NormalTok{    "x": \{"field": "x", "type": "quantitative"\},}
\NormalTok{    "y": \{"field": "y", "type": "quantitative"\}}
\NormalTok{  \}}
\NormalTok{\}}
\end{Highlighting}
\end{Shaded}

\subsubsection{More chart types}\label{more-chart-types}

The four shapes the chart builder can also produce: a histogram, which
bins one column and counts the rows in each bin; a heatmap, which
colours a grid by a quantity; error bars, which summarise repeated
measurements; and a pie, whose slice size is \texttt{theta} and whose
category is the colour. Each carries a \texttt{title}, which --- as
section 8 explains --- makes it a numbered figure.

\paragraph{Markdown}\label{markdown-7}

\begin{Shaded}
\begin{Highlighting}[]
\InformationTok{\textasciigrave{}\textasciigrave{}\textasciigrave{}vega{-}lite}
\InformationTok{\{}
\InformationTok{  "title": "Distribution of recovered masses",}
\InformationTok{  "data": \{"values": [}
\InformationTok{    \{"mass": 1.2\}, \{"mass": 1.9\}, \{"mass": 2.4\}, \{"mass": 2.6\}, \{"mass": 3.1\}}
\InformationTok{  ]\},}
\InformationTok{  "mark": "bar",}
\InformationTok{  "encoding": \{}
\InformationTok{    "x": \{"field": "mass", "bin": true, "type": "quantitative"\},}
\InformationTok{    "y": \{"aggregate": "count", "type": "quantitative"\}}
\InformationTok{  \}}
\InformationTok{\}}
\InformationTok{\textasciigrave{}\textasciigrave{}\textasciigrave{}}

\InformationTok{\textasciigrave{}\textasciigrave{}\textasciigrave{}vega{-}lite}
\InformationTok{\{}
\InformationTok{  "title": "Detections by day and hour",}
\InformationTok{  "data": \{"values": [}
\InformationTok{    \{"day": "Mon", "hour": 9, "rate": 4\}, \{"day": "Mon", "hour": 10, "rate": 7\},}
\InformationTok{    \{"day": "Tue", "hour": 9, "rate": 2\}, \{"day": "Tue", "hour": 10, "rate": 9\}}
\InformationTok{  ]\},}
\InformationTok{  "mark": "rect",}
\InformationTok{  "encoding": \{}
\InformationTok{    "x": \{"field": "day", "type": "nominal"\},}
\InformationTok{    "y": \{"field": "hour", "type": "nominal"\},}
\InformationTok{    "color": \{"field": "rate", "type": "quantitative", "aggregate": "mean"\}}
\InformationTok{  \}}
\InformationTok{\}}
\InformationTok{\textasciigrave{}\textasciigrave{}\textasciigrave{}}

\InformationTok{\textasciigrave{}\textasciigrave{}\textasciigrave{}vega{-}lite}
\InformationTok{\{}
\InformationTok{  "title": "Yield by variety, with 95\% confidence intervals",}
\InformationTok{  "data": \{"values": [}
\InformationTok{    \{"variety": "A", "yield": 21\}, \{"variety": "A", "yield": 25\},}
\InformationTok{    \{"variety": "A", "yield": 23\}, \{"variety": "B", "yield": 31\},}
\InformationTok{    \{"variety": "B", "yield": 28\}, \{"variety": "B", "yield": 34\}}
\InformationTok{  ]\},}
\InformationTok{  "mark": \{"type": "errorbar", "extent": "ci"\},}
\InformationTok{  "encoding": \{}
\InformationTok{    "x": \{"field": "variety", "type": "nominal"\},}
\InformationTok{    "y": \{"field": "yield", "type": "quantitative"\}}
\InformationTok{  \}}
\InformationTok{\}}
\InformationTok{\textasciigrave{}\textasciigrave{}\textasciigrave{}}

\InformationTok{\textasciigrave{}\textasciigrave{}\textasciigrave{}vega{-}lite}
\InformationTok{\{}
\InformationTok{  "title": "Share of detections by instrument",}
\InformationTok{  "data": \{"values": [}
\InformationTok{    \{"instrument": "A", "count": 12\}, \{"instrument": "B", "count": 7\},}
\InformationTok{    \{"instrument": "C", "count": 4\}}
\InformationTok{  ]\},}
\InformationTok{  "mark": "arc",}
\InformationTok{  "encoding": \{}
\InformationTok{    "theta": \{"field": "count", "type": "quantitative"\},}
\InformationTok{    "color": \{"field": "instrument", "type": "nominal"\}}
\InformationTok{  \}}
\InformationTok{\}}
\InformationTok{\textasciigrave{}\textasciigrave{}\textasciigrave{}}
\end{Highlighting}
\end{Shaded}

\paragraph{Rendered}\label{rendered-7}

\begin{Shaded}
\begin{Highlighting}[]
\NormalTok{\{}
\NormalTok{  "title": "Distribution of recovered masses",}
\NormalTok{  "data": \{"values": [}
\NormalTok{    \{"mass": 1.2\}, \{"mass": 1.9\}, \{"mass": 2.4\}, \{"mass": 2.6\}, \{"mass": 3.1\}}
\NormalTok{  ]\},}
\NormalTok{  "mark": "bar",}
\NormalTok{  "encoding": \{}
\NormalTok{    "x": \{"field": "mass", "bin": true, "type": "quantitative"\},}
\NormalTok{    "y": \{"aggregate": "count", "type": "quantitative"\}}
\NormalTok{  \}}
\NormalTok{\}}
\end{Highlighting}
\end{Shaded}

\begin{Shaded}
\begin{Highlighting}[]
\NormalTok{\{}
\NormalTok{  "title": "Detections by day and hour",}
\NormalTok{  "data": \{"values": [}
\NormalTok{    \{"day": "Mon", "hour": 9, "rate": 4\}, \{"day": "Mon", "hour": 10, "rate": 7\},}
\NormalTok{    \{"day": "Tue", "hour": 9, "rate": 2\}, \{"day": "Tue", "hour": 10, "rate": 9\}}
\NormalTok{  ]\},}
\NormalTok{  "mark": "rect",}
\NormalTok{  "encoding": \{}
\NormalTok{    "x": \{"field": "day", "type": "nominal"\},}
\NormalTok{    "y": \{"field": "hour", "type": "nominal"\},}
\NormalTok{    "color": \{"field": "rate", "type": "quantitative", "aggregate": "mean"\}}
\NormalTok{  \}}
\NormalTok{\}}
\end{Highlighting}
\end{Shaded}

\begin{Shaded}
\begin{Highlighting}[]
\NormalTok{\{}
\NormalTok{  "title": "Yield by variety, with 95\% confidence intervals",}
\NormalTok{  "data": \{"values": [}
\NormalTok{    \{"variety": "A", "yield": 21\}, \{"variety": "A", "yield": 25\},}
\NormalTok{    \{"variety": "A", "yield": 23\}, \{"variety": "B", "yield": 31\},}
\NormalTok{    \{"variety": "B", "yield": 28\}, \{"variety": "B", "yield": 34\}}
\NormalTok{  ]\},}
\NormalTok{  "mark": \{"type": "errorbar", "extent": "ci"\},}
\NormalTok{  "encoding": \{}
\NormalTok{    "x": \{"field": "variety", "type": "nominal"\},}
\NormalTok{    "y": \{"field": "yield", "type": "quantitative"\}}
\NormalTok{  \}}
\NormalTok{\}}
\end{Highlighting}
\end{Shaded}

\begin{Shaded}
\begin{Highlighting}[]
\NormalTok{\{}
\NormalTok{  "title": "Share of detections by instrument",}
\NormalTok{  "data": \{"values": [}
\NormalTok{    \{"instrument": "A", "count": 12\}, \{"instrument": "B", "count": 7\},}
\NormalTok{    \{"instrument": "C", "count": 4\}}
\NormalTok{  ]\},}
\NormalTok{  "mark": "arc",}
\NormalTok{  "encoding": \{}
\NormalTok{    "theta": \{"field": "count", "type": "quantitative"\},}
\NormalTok{    "color": \{"field": "instrument", "type": "nominal"\}}
\NormalTok{  \}}
\NormalTok{\}}
\end{Highlighting}
\end{Shaded}

\subsection{8. Figures}\label{figures}

A caption is what makes a figure. A chart's \texttt{title}, an image's
alt text or a Mermaid diagram's \texttt{title} becomes the caption,
drawn once below the figure, and every captioned figure is numbered in
document order --- so the four titled charts above are Figures 1--4 and
the sequence continues here. An image with empty alt text, or an
untitled diagram, stays decorative: no number, no caption. Images may be
remote URLs or paths relative to the document.

\subsubsection{Markdown}\label{markdown-8}

\begin{Shaded}
\begin{Highlighting}[]
\InformationTok{\textasciigrave{}\textasciigrave{}\textasciigrave{}vega{-}lite}
\InformationTok{\{}
\InformationTok{  "title": "Recovered sources",}
\InformationTok{  "data": \{"values": [}
\InformationTok{    \{"category": "A", "value": 12\}, \{"category": "B", "value": 27\},}
\InformationTok{    \{"category": "C", "value": 19\}}
\InformationTok{  ]\},}
\InformationTok{  "mark": "bar",}
\InformationTok{  "encoding": \{}
\InformationTok{    "x": \{"field": "category", "type": "nominal"\},}
\InformationTok{    "y": \{"field": "value", "type": "quantitative"\}}
\InformationTok{  \}}
\InformationTok{\}}
\InformationTok{\textasciigrave{}\textasciigrave{}\textasciigrave{}}

\NormalTok{A captioned remote image:}

\AlertTok{![A placeholder photograph](https://placehold.co/400x250)}

\NormalTok{A decorative image, with empty alt text:}

\AlertTok{![](https://placehold.co/120x80)}

\NormalTok{A local image, resolved relative to this document\textquotesingle{}s own folder:}

\AlertTok{![A locally stored figure](sample{-}figure.png)}

\NormalTok{A chart that sets its own }\InformationTok{\textasciigrave{}"width"\textasciigrave{}}\NormalTok{, overriding View → Chart Width:}

\InformationTok{\textasciigrave{}\textasciigrave{}\textasciigrave{}vega{-}lite}
\InformationTok{\{}
\InformationTok{  "width": 150,}
\InformationTok{  "data": \{"values": [\{"x": 0, "y": 1\}, \{"x": 1, "y": 2\}, \{"x": 2, "y": 4\}]\},}
\InformationTok{  "mark": "point",}
\InformationTok{  "encoding": \{}
\InformationTok{    "x": \{"field": "x", "type": "quantitative"\},}
\InformationTok{    "y": \{"field": "y", "type": "quantitative"\}}
\InformationTok{  \}}
\InformationTok{\}}
\InformationTok{\textasciigrave{}\textasciigrave{}\textasciigrave{}}

\NormalTok{A titled Mermaid diagram:}

\InformationTok{\textasciigrave{}\textasciigrave{}\textasciigrave{}mermaid}
\InformationTok{{-}{-}{-}}
\InformationTok{title: Pipeline stages}
\InformationTok{{-}{-}{-}}
\InformationTok{flowchart LR}
\InformationTok{  A[Ingest] {-}{-}\textgreater{} B[Parse]}
\InformationTok{  B {-}{-}\textgreater{} C[Render]}
\InformationTok{\textasciigrave{}\textasciigrave{}\textasciigrave{}}

\NormalTok{An untitled diagram of a different kind:}

\InformationTok{\textasciigrave{}\textasciigrave{}\textasciigrave{}mermaid}
\InformationTok{sequenceDiagram}
\InformationTok{  participant Author}
\InformationTok{  participant Halmos}
\InformationTok{  Author{-}\textgreater{}\textgreater{}Halmos: Type a mermaid fence}
\InformationTok{  Halmos{-}{-}\textgreater{}\textgreater{}Author: Render diagram}
\InformationTok{\textasciigrave{}\textasciigrave{}\textasciigrave{}}
\end{Highlighting}
\end{Shaded}

\subsubsection{Rendered}\label{rendered-8}

\begin{Shaded}
\begin{Highlighting}[]
\NormalTok{\{}
\NormalTok{  "title": "Recovered sources",}
\NormalTok{  "data": \{"values": [}
\NormalTok{    \{"category": "A", "value": 12\}, \{"category": "B", "value": 27\},}
\NormalTok{    \{"category": "C", "value": 19\}}
\NormalTok{  ]\},}
\NormalTok{  "mark": "bar",}
\NormalTok{  "encoding": \{}
\NormalTok{    "x": \{"field": "category", "type": "nominal"\},}
\NormalTok{    "y": \{"field": "value", "type": "quantitative"\}}
\NormalTok{  \}}
\NormalTok{\}}
\end{Highlighting}
\end{Shaded}

A captioned remote image:

\begin{figure}
\centering
\pandocbounded{\includegraphics[keepaspectratio,alt={A placeholder photograph}]{https://placehold.co/400x250}}
\caption{A placeholder photograph}
\end{figure}

A decorative image, with empty alt text:

\pandocbounded{\includegraphics[keepaspectratio]{https://placehold.co/120x80}}

A local image, resolved relative to this document's own folder:

\begin{figure}
\centering
\pandocbounded{\includegraphics[keepaspectratio,alt={A locally stored figure}]{sample-figure.png}}
\caption{A locally stored figure}
\end{figure}

A chart that sets its own \texttt{"width"}, overriding View → Chart
Width:

\begin{Shaded}
\begin{Highlighting}[]
\NormalTok{\{}
\NormalTok{  "width": 150,}
\NormalTok{  "data": \{"values": [\{"x": 0, "y": 1\}, \{"x": 1, "y": 2\}, \{"x": 2, "y": 4\}]\},}
\NormalTok{  "mark": "point",}
\NormalTok{  "encoding": \{}
\NormalTok{    "x": \{"field": "x", "type": "quantitative"\},}
\NormalTok{    "y": \{"field": "y", "type": "quantitative"\}}
\NormalTok{  \}}
\NormalTok{\}}
\end{Highlighting}
\end{Shaded}

A titled Mermaid diagram:

\begin{Shaded}
\begin{Highlighting}[]
\NormalTok{{-}{-}{-}}
\NormalTok{title: Pipeline stages}
\NormalTok{{-}{-}{-}}
\NormalTok{flowchart LR}
\NormalTok{  A[Ingest] {-}{-}\textgreater{} B[Parse]}
\NormalTok{  B {-}{-}\textgreater{} C[Render]}
\end{Highlighting}
\end{Shaded}

An untitled diagram of a different kind:

\begin{Shaded}
\begin{Highlighting}[]
\NormalTok{sequenceDiagram}
\NormalTok{  participant Author}
\NormalTok{  participant Halmos}
\NormalTok{  Author{-}\textgreater{}\textgreater{}Halmos: Type a mermaid fence}
\NormalTok{  Halmos{-}{-}\textgreater{}\textgreater{}Author: Render diagram}
\end{Highlighting}
\end{Shaded}

View → Figure Alignment moves every figure and decorative image
together, captions included; this PDF shows whichever alignment was set
when it was exported.

\subsection{9. Citations}\label{citations}

Pandoc-style keys --- \texttt{{[}@key{]}}, \texttt{@key} for a narrative
citation, \texttt{-@key} to suppress the author, and a locator after a
comma --- are resolved against the \texttt{.bib} named in the
frontmatter and formatted in the chosen CSL style. Every work cited
collects into the References list at the end of the document.

\subsubsection{Markdown}\label{markdown-9}

\begin{Shaded}
\begin{Highlighting}[]
\NormalTok{A simple citation }\CommentTok{[}\OtherTok{@zurekQuantumTheoryClassical2018}\CommentTok{]}\NormalTok{, a multi{-}cite}
\CommentTok{[}\OtherTok{@everettRelativeStateFormulation1957; @zurekQuantumTheoryClassical2018}\CommentTok{]}\NormalTok{, and a}
\NormalTok{narrative one: @korbiczRoadsObjectivityQuantum2021 sets out the argument.}

\NormalTok{Suppressed author, because the year is already in the sentence — Everett\textquotesingle{}s 1957}
\NormalTok{paper }\CommentTok{[}\OtherTok{{-}@everettRelativeStateFormulation1957}\CommentTok{]}\NormalTok{ — and a locator attached to a key}
\CommentTok{[}\OtherTok{see @kloeffelProspectsSpinBasedQuantum2013, pp. 51{-}53}\CommentTok{]}\NormalTok{. A section locator}
\NormalTok{works the same way }\CommentTok{[}\OtherTok{@hanceWhatDoesIt2022, sec. 2}\CommentTok{]}\NormalTok{.}

\NormalTok{A group split by a hard line wrap formats like any other: [see}
\NormalTok{@leStrongQuantumDarwinism2019, pp. 3{-}4].}

\NormalTok{Several works supporting one claim [@horodeckiQuantumOriginsObjectivity2015;}
\NormalTok{@undenRevealingEmergenceClassicality2019; @korbiczRoadsObjectivityQuantum2021],}
\NormalTok{and a citation carried into a quotation intact:}

\AttributeTok{\textgreater{} Objectivity is emergent }\CommentTok{[}\OtherTok{@leStrongQuantumDarwinism2019}\CommentTok{]}\AttributeTok{.}

\NormalTok{Eight authors collapse to *et al.* }\CommentTok{[}\OtherTok{@feinQuantumSuperpositionMolecules2019}\CommentTok{]}\NormalTok{. A}
\NormalTok{work with no year reads *n.d.* }\CommentTok{[}\OtherTok{@fineBohrsResponseEPR}\CommentTok{]}\NormalTok{. A work with no author}
\NormalTok{falls back to its title }\CommentTok{[}\OtherTok{@DecoherenceQuantumToClassicalTransition2007}\CommentTok{]}\NormalTok{. A}
\NormalTok{preprint renders as one }\CommentTok{[}\OtherTok{@mullerSixMeasurementProblems2023}\CommentTok{]}\NormalTok{.}

\NormalTok{An unknown key renders as a visible inline error without blanking the rest of}
\NormalTok{the page: }\CommentTok{[}\OtherTok{@notakey1999}\CommentTok{]}\NormalTok{. A group mixing a good key with a bad one is reported}
\NormalTok{the same way }\CommentTok{[}\OtherTok{@zurekQuantumTheoryClassical2018; @alsomissing1999}\CommentTok{]}\NormalTok{.}
\end{Highlighting}
\end{Shaded}

\subsubsection{Rendered}\label{rendered-9}

A simple citation (Zurek, 2018), a multi-cite (Everett, 1957; Zurek,
2018), and a narrative one: Korbicz (2021) sets out the argument.

Suppressed author, because the year is already in the sentence ---
Everett's 1957 paper (1957) --- and a locator attached to a key (see
Kloeffel \& Loss, 2013, pp. 51--53). A section locator works the same
way (Hance \& Hossenfelder, 2022, Section 2).

A group split by a hard line wrap formats like any other: (see Le \&
Olaya-Castro, 2019, pp. 3--4).

Several works supporting one claim (Horodecki et al., 2015; Korbicz,
2021; Unden et al., 2019), and a citation carried into a quotation
intact:

\begin{quote}
Objectivity is emergent (Le \& Olaya-Castro, 2019).
\end{quote}

Eight authors collapse to \emph{et al.} (Fein et al., 2019). A work with
no year reads \emph{n.d.} (Fine, n.d.). A work with no author falls back
to its title (\emph{Decoherence and the {Quantum-To-Classical
Transition}}, 2007). A preprint renders as one (Muller, 2023).

An unknown key renders as a visible inline error without blanking the
rest of the page: (\textbf{notakey1999?}). A group mixing a good key
with a bad one is reported the same way (Zurek, 2018;
\textbf{alsomissing1999?}).

Titles in this export carry BibTeX capital-protection like
\texttt{\{\{Quantum\ Darwinism\}\}}, and surnames like
\texttt{\{Olaya-Castro\}} are braced to keep them whole; the References
list shows \emph{Quantum Darwinism} and \emph{Olaya-Castro}, never the
braces. Author-year disambiguation (\texttt{2021a} / \texttt{2021b})
cannot be shown here, because no two works in this library share a first
author and a year.

\subsection{10. Things that do not become
citations}\label{things-that-do-not-become-citations}

None of the following turns into a citation or reaches the References
list.

\subsubsection{Markdown}\label{markdown-10}

\begin{Shaded}
\begin{Highlighting}[]
\SpecialStringTok{{-} }\NormalTok{An email address: write to a.zurek@example.org for the dataset.}
\SpecialStringTok{{-} }\NormalTok{Plain brackets used as brackets: }\CommentTok{[}\OtherTok{see the appendix}\CommentTok{]}\NormalTok{ and }\CommentTok{[}\OtherTok{2}\CommentTok{]}\NormalTok{.}
\SpecialStringTok{{-} }\NormalTok{A markdown link: }\CommentTok{[}\OtherTok{the project page}\CommentTok{](https://example.org/project)}\NormalTok{.}
\SpecialStringTok{{-} }\NormalTok{A citation inside a code span: }\InformationTok{\textasciigrave{}[@zurekQuantumTheoryClassical2018]\textasciigrave{}}\NormalTok{.}
\SpecialStringTok{{-} }\NormalTok{A fenced code block:}

\InformationTok{\textasciigrave{}\textasciigrave{}\textasciigrave{}text}
\InformationTok{[@everettRelativeStateFormulation1957] is not formatted here either.}
\InformationTok{\textasciigrave{}\textasciigrave{}\textasciigrave{}}
\end{Highlighting}
\end{Shaded}

\subsubsection{Rendered}\label{rendered-10}

\begin{itemize}
\tightlist
\item
  An email address: write to a.zurek@example.org for the dataset.
\item
  Plain brackets used as brackets: {[}see the appendix{]} and {[}2{]}.
\item
  A markdown link: \href{https://example.org/project}{the project page}.
\item
  A citation inside a code span:
  \texttt{{[}@zurekQuantumTheoryClassical2018{]}}.
\item
  A fenced code block:
\end{itemize}

\begin{Shaded}
\begin{Highlighting}[]
\NormalTok{[@everettRelativeStateFormulation1957] is not formatted here either.}
\end{Highlighting}
\end{Shaded}

\subsection{11. Errors are contained}\label{errors-are-contained}

Mistakes render in place, and everything after them still renders. These
three are \textbf{supposed} to look broken.

\subsubsection{Markdown}\label{markdown-11}

\begin{Shaded}
\begin{Highlighting}[]
\AnnotationTok{Invalid LaTeX renders inline in red:}\CommentTok{ $\textbackslash{}thisisnotacommand\{x\}$}

\NormalTok{Broken Vega{-}Lite JSON renders an error card:}

\InformationTok{\textasciigrave{}\textasciigrave{}\textasciigrave{}vega{-}lite}
\InformationTok{\{ "mark": "bar", this is not valid JSON}
\InformationTok{\textasciigrave{}\textasciigrave{}\textasciigrave{}}

\NormalTok{Broken Mermaid syntax renders the same kind of card:}

\InformationTok{\textasciigrave{}\textasciigrave{}\textasciigrave{}mermaid}
\InformationTok{this is not a valid diagram type}
\InformationTok{\textasciigrave{}\textasciigrave{}\textasciigrave{}}
\end{Highlighting}
\end{Shaded}

\subsubsection{Rendered}\label{rendered-11}

Invalid LaTeX renders inline in red: \(\thisisnotacommand{x}\)

Broken Vega-Lite JSON renders an error card:

\begin{Shaded}
\begin{Highlighting}[]
\NormalTok{\{ "mark": "bar", this is not valid JSON}
\end{Highlighting}
\end{Shaded}

Broken Mermaid syntax renders the same kind of card:

\begin{Shaded}
\begin{Highlighting}[]
\NormalTok{this is not a valid diagram type}
\end{Highlighting}
\end{Shaded}

If you can read this line with normal styling, error containment works.
✅

\subsection{12. Beyond the page}\label{beyond-the-page}

Some of Halmos lives in the editor rather than on the page, and a PDF
cannot show it. Briefly, for completeness:

\begin{itemize}
\tightlist
\item
  \textbf{Chart builder.} Insert → Chart\ldots{} builds a
  \texttt{vega-lite} block from a pasted or imported table, and reopens
  an existing block with every control --- data included --- prefilled.
  A chart using a property the builder cannot express, like the
  \texttt{axis} in section 7's bar chart, is refused with an
  explanation.
\item
  \textbf{Formatting commands.} ⌘B, ⌘I and ⌘⇧K toggle bold, italic and
  code; ⌘1--⌘6 set headings; ⌘⇧8 and ⌘⇧7 make bullet and numbered lists.
  They never rewrite frontmatter, fenced code or inline code spans.
\item
  \textbf{Folding.} ⌘⌥{[} and ⌘⌥{]} fold and unfold the block at the
  cursor; View → Fold All Code Blocks collapses every fence to its
  opening line.
\item
  \textbf{Live bibliography.} Saving the \texttt{.bib} from Zotero while
  the document is open refreshes the preview on its own.
\item
  \textbf{Appearance.} View → Appearance follows System, Light or Dark;
  the exported PDF is always light, whatever the screen theme, with
  figures kept whole across page breaks.
\end{itemize}

\section{Test copy an image}\label{test-copy-an-image}

\begin{figure}
\centering
\pandocbounded{\includegraphics[keepaspectratio,alt={A Bloch Sphere}]{images/pasted-20261018-130949.png}}
\caption{A Bloch Sphere}
\end{figure}

\protect\phantomsection\label{refs}
\begin{CSLReferences}{1}{0}
\bibitem[\citeproctext]{ref-DecoherenceQuantumToClassicalTransition2007}
\emph{Decoherence and the {Quantum-To-Classical Transition}}. (2007).
Springer Berlin Heidelberg.
\url{https://doi.org/10.1007/978-3-540-35775-9}

\bibitem[\citeproctext]{ref-everettRelativeStateFormulation1957}
Everett, H. (1957). "{Relative State}" {Formulation} of {Quantum
Mechanics}. \emph{Reviews of Modern Physics}, \emph{29}(3), 454--462.
\url{https://doi.org/10.1103/RevModPhys.29.454}

\bibitem[\citeproctext]{ref-feinQuantumSuperpositionMolecules2019}
Fein, Y. Y., Geyer, P., Zwick, P., Kiałka, F., Pedalino, S., Mayor, M.,
Gerlich, S., \& Arndt, M. (2019). Quantum superposition of molecules
beyond 25 {kDa}. \emph{Nature Physics}, \emph{15}(12), 1242--1245.
\url{https://doi.org/10.1038/s41567-019-0663-9}

\bibitem[\citeproctext]{ref-fineBohrsResponseEPR}
Fine, A. (n.d.). \emph{Bohr's {Response} to {EPR}: {Criticism} and
{Defense}}.

\bibitem[\citeproctext]{ref-hanceWhatDoesIt2022}
Hance, J. R., \& Hossenfelder, S. (2022). What does it take to solve the
measurement problem? \emph{Journal of Physics Communications},
\emph{6}(10), 102001. \url{https://doi.org/10.1088/2399-6528/ac96cf}

\bibitem[\citeproctext]{ref-horodeckiQuantumOriginsObjectivity2015}
Horodecki, R., Korbicz, J. K., \& Horodecki, P. (2015). Quantum origins
of objectivity. \emph{Physical Review A}, \emph{91}(3), 032122.
\url{https://doi.org/10.1103/PhysRevA.91.032122}

\bibitem[\citeproctext]{ref-kloeffelProspectsSpinBasedQuantum2013}
Kloeffel, C., \& Loss, D. (2013). Prospects for {Spin-Based Quantum
Computing} in {Quantum Dots}. \emph{Annual Review of Condensed Matter
Physics}, \emph{4}(1), 51--81.
\url{https://doi.org/10.1146/annurev-conmatphys-030212-184248}

\bibitem[\citeproctext]{ref-korbiczRoadsObjectivityQuantum2021}
Korbicz, J. K. (2021). Roads to objectivity: {Quantum Darwinism},
{Spectrum Broadcast Structures}, and {Strong} quantum {Darwinism} -- a
review. \emph{Quantum}, \emph{5}, 571.
\url{https://doi.org/10.22331/q-2021-11-08-571}

\bibitem[\citeproctext]{ref-leStrongQuantumDarwinism2019}
Le, T. P., \& Olaya-Castro, A. (2019). Strong {Quantum Darwinism} and
{Strong Independence} are {Equivalent} to {Spectrum Broadcast
Structure}. \emph{Physical Review Letters}, \emph{122}(1), 010403.
\url{https://doi.org/10.1103/PhysRevLett.122.010403}

\bibitem[\citeproctext]{ref-mullerSixMeasurementProblems2023}
Muller, F. A. (2023). \emph{Six {Measurement Problems} of {Quantum
Mechanics}} (arXiv:2305.10206). arXiv.
\url{https://doi.org/10.48550/arXiv.2305.10206}

\bibitem[\citeproctext]{ref-undenRevealingEmergenceClassicality2019}
Unden, T. K., Louzon, D., Zwolak, M., Zurek, W. H., \& Jelezko, F.
(2019). Revealing the {Emergence} of {Classicality Using
Nitrogen-Vacancy Centers}. \emph{Physical Review Letters},
\emph{123}(14), 140402.
\url{https://doi.org/10.1103/PhysRevLett.123.140402}

\bibitem[\citeproctext]{ref-zurekQuantumTheoryClassical2018}
Zurek, W. H. (2018). Quantum theory of the classical: Quantum jumps,
{Born}'s {Rule} and objective classical reality via quantum {Darwinism}.
\emph{Philosophical Transactions of the Royal Society A: Mathematical,
Physical and Engineering Sciences}, \emph{376}(2123), 20180107.
\url{https://doi.org/10.1098/rsta.2018.0107}

\end{CSLReferences}

\end{document}
